\documentclass[runningheads]{llncs}

\usepackage[T1]{fontenc}
\usepackage{graphicx}
\usepackage{booktabs}
\usepackage{amsmath}
\usepackage{amssymb}
\providecommand{\mathbbm}[1]{\mathbf{#1}}
\usepackage{enumitem}
\usepackage{multirow}
\usepackage{xcolor}
\usepackage[colorlinks=true,linkcolor=blue,urlcolor=blue,citecolor=blue]{hyperref}

\begin{document}

\title{Pref-Rerank: User-Conditioned Bilinear Score Intervention for Personalized Evidence Reranking}

\titlerunning{Pref-Rerank for Personalized Web-RAG}

\author{Anonymous Author(s)}
\authorrunning{Anonymous}
\institute{Anonymous Affiliation\\
\email{anonymous@example.com}}

\maketitle

\begin{abstract}
Query-log-based evidence reranking systems serve every user the same top-$k$ evidence passages, ignoring the evidence-type preferences that are implicit in a user's search history. We propose \emph{Pref-Rerank}, a minimal user-conditioning mechanism for the reranking stage: a single bilinear user$\times$document interaction term $h(q,d)^\top A\, u$ added to a frozen scorer's output, where $u \in \mathbb{R}^{16}$ is a compact user vector derived from session history. The module has only \textbf{12{,}288 trainable parameters} and trains in seconds on top of precomputed frozen features. On a 200k-session sample of the AOL query log ($n{=}201$ test sessions), the behavioral-features variant of Pref-Rerank improves session-level nDCG@10 from $0.459$ (user-agnostic) to $0.543$ ($\Delta{=}{+}0.084$, paired bootstrap 95\% CI $[+0.034, +0.131]$). The LLM-extracted typed-schema variant is also significant but weaker ($\Delta{=}{+}0.054$, CI $[+0.013, +0.094]$). Two findings stand out. First, the bilinear interaction $h^\top A u$ is the active mechanism: a user-only bias $w^\top u$ without document interaction contributes nothing. Second, random user vectors do \emph{not} yield a significant gain ($\Delta{=}{+}0.027$, CI $[-0.010, +0.065]$), which rules out the concern that the bilinear term merely memorizes session identity rather than modeling user signal. Behavioral features provide the strongest and most stable user representation; LLM-extracted schemas retain an interpretability advantage but not a ranking-quality advantage.

\keywords{Web information retrieval \and Retrieval-augmented generation \and Personalized search \and Cross-encoder reranking \and Click-log learning}
\end{abstract}

\section{Introduction}
\label{sec:intro}

Web-backed retrieval-augmented generation (RAG) systems~\cite{lewis2020rag,sun2024ragsurvey} answer user queries by retrieving web pages and conditioning a language model on the retrieved evidence. To keep latency and fetch cost low, these systems cache retrieved pages and re-rank a candidate set before generation. The re-ranking stage typically uses a strong frozen scorer such as a cross-encoder~\cite{nogueira2020monot5,qin2023llmrerank} that scores each (query, document) pair independently. A limitation of this design is that the score depends only on the query and the document: every user issuing the same query receives the same ranked evidence.

This matters because the same query can mean different things to different users. A student searching ``MacBook'' and a designer searching ``MacBook'' may get the same candidate pages from the retriever, but one wants price comparisons and the other wants display specs. The cross-encoder has no way to tell them apart. The signal to disambiguate is in the session history, but the reranking stage ignores it.

Existing approaches to personalized search fall into three families, each with a drawback. Click-log learning-to-rank~\cite{craswell2007click,joachims2007evaluating} models users through behavioral features but carries no semantic preference structure and is opaque. Prompt-based personalization conditions the reranker or generator on the raw history text, which is fragile under cold start and offers no learnable user representation. Heavier personalized reranker heads, which append learnable user modules to a cross-encoder, work but obscure the question we want to answer: \emph{what simple mechanism can make a frozen scorer user-aware?}

We answer this question with \emph{Pref-Rerank}. Given a frozen scorer producing a base score $s_{\mathrm{base}}(q,d)$ and a document representation $h(q,d)\in\mathbb{R}^{768}$, we add one bilinear term $h(q,d)^\top A\, u$, where $u\in\mathbb{R}^{16}$ is a user vector and $A\in\mathbb{R}^{768\times 16}$ is the only trainable component (12{,}288 parameters). We study three ways to build $u$: an LLM-extracted typed preference schema, a six-dimensional behavioral feature vector, and a random vector as control. Training uses a session-level pairwise ranking loss on position-debiased (clicked, not-clicked) pairs from the AOL query log~\cite{pass2006_picture}.

Two findings stand out. First, the bilinear term works: on a 200k-session sample of AOL (of which 201 test sessions survive candidate-set filtering) ($n{=}201$ test sessions), the behavioral-features variant raises session-level nDCG@10 from $0.459$ to $0.543$ ($\Delta{=}{+}0.084$, paired bootstrap 95\% CI $[+0.034, +0.131]$), with 12{,}288 trainable parameters. The LLM-schema variant is also significant but weaker ($\Delta{=}{+}0.054$). Second, a random user vector does \emph{not} yield a significant gain ($\Delta{=}{+}0.027$, CI $[-0.010, +0.065]$) --- the bilinear term exploits real user signal, not merely session identity.. A user-only bias $w^\top u$ without document interaction also contributes nothing. Figure~\ref{fig:hero} shows both results.

\begin{figure}[t]
\centering
\includegraphics[width=\linewidth]{figures/fig1_hero_ndcg.pdf}
\caption{Session-level nDCG@10 across six systems on 200k AOL sessions ($n{=}201$, error bars: paired bootstrap 95\% CI). Bilinear variants with real user signal (\texttt{classic} behavioral, \texttt{no\_user\_bias} LLM schema) significantly beat the baseline; \texttt{random\_schema} and \texttt{linear\_only} ($w^\top u$ only) do not. Behavioral features are the strongest user representation; LLM schema retains interpretability but not a ranking-quality edge.}
\label{fig:hero}
\end{figure}

What this leaves open, and what the data argue against, is whether the bilinear gain depends on the specific user representation. The random-vector control (no signal, just a fixed per-session code) produces only $+0.027$, a non-significant difference; LLM-extracted typed schemas ($+0.054$, significant) and behavioral features ($+0.084$, significant) outperform it. The user vector is not the source of the gain --- the bilinear interaction mechanism is. We discuss in Sec.~\ref{sec:discussion} the possibility that the bilinear term partly models session identity rather than preference, and the limits of what $n{=}201$ can establish.

\paragraph{Contributions.}
\begin{itemize}[leftmargin=*,itemsep=2pt]
\item A simple effective bilinear score intervention for frozen scorers: one trainable matrix ($12{,}288$ parameters) improves session-level nDCG@10 by $+0.084$ on 200k AOL sessions ($n{=}201$, paired bootstrap 95\% CI $[+0.034, +0.131]$, Sec.~\ref{sec:exp}).
\item A mechanism ablation isolating the bilinear term $h^\top A u$ as the active component: the user-only bias $w^\top u$ alone is indistinguishable from baseline (Sec.~\ref{sec:exp:mechanism}).
\item A user-representation ablation with a clean control: behavioral features ($+0.084$) and LLM-extracted typed schemas ($+0.054$) are both significant, while random user vectors are not ($+0.027$, CI crosses zero). This does not support the session-identity-memorization hypothesis and shows the bilinear term exploits genuine user signal (Sec.~\ref{sec:exp:representation}).
\end{itemize}

\section{Related Work}
\label{sec:related}

\paragraph{Retrieval-augmented generation.}
RAG~\cite{lewis2020rag} has become the standard paradigm for grounding LLM generation in external knowledge. Recent work extends it with self-reflection~\cite{asai2024selfrag}, active retrieval~\cite{jiang2023adaptive}, and structured retrieval planning~\cite{gao2023selfrag}. Comprehensive surveys~\cite{sun2024ragsurvey} systematize the growing design space. Our work does not modify the retrieval or generation stages; it targets the evidence \emph{reranking} stage that sits between them, and shows that a compact user-conditioning term on a frozen reranker can improve what evidence reaches the generator.

\paragraph{Personalized search and click-log learning-to-rank.}
A long line of work models the user from search logs. Click-bipartite and session-based approaches~\cite{craswell2007click,joachims2017unbiased} treat clicks as implicit feedback and learn ranking functions from click logs, with explicit attention to position bias~\cite{joachims2007evaluating,craswell2008experimental}. These methods are practical and deployable but model the user through behavioral aggregates: click entropy, dwell-time proxies, reformulation patterns --- none of which carry semantic preference structure or support inspection. Pref-Rerank uses the same AOL click signal for supervision but introduces a structured, learnable user vector and intervenes at the rerank score rather than in the feature space of a learning-to-rank model.

\paragraph{LLM-based reranking and user preference extraction.}
Large language models have been used directly for passage reranking through pairwise~\cite{qin2023llmrerank,gao2024pairwise} and listwise~\cite{pradeep2023rankgpt} prompting. PMG~\cite{shi2024pmg} extracts both explicit preference keywords and implicit soft embeddings from user history for personalized multimodal generation. Our schema-constrained extraction follows a similar philosophy but targets the \emph{reranking} stage, and our ablation (Sec.~\ref{sec:exp}) shows that the specific extraction is not the source of the gain --- the bilinear injection mechanism is. Structured context extraction for personalization~\cite{zhu2024personalized} is a related direction; we differ by injecting the extracted structure into a learnable bilinear score rather than a prompt.

\paragraph{Cross-encoder and dense retrieval.}
Cross-encoders such as monoT5~\cite{nogueira2020monot5} and late-interaction models such as ColBERT~\cite{khattab2020colbert} set the state of the art on passage retrieval and reranking. Dense passage retrieval~\cite{karpukhin2020dpr} and frozen encoders like BGE~\cite{xiao2023bge} provide strong document representations. We do not propose a new cross-encoder; we propose a lightweight intervention applicable to any frozen scorer that exposes a scalar score and a document representation, instantiated on a BM25~\cite{robertson2009bm25} + BGE query-flow backbone suited to the click-log setting.

\paragraph{Recommendation-style bilinear user modeling.}
The bilinear form $h^\top A u$ is classical in recommendation~\cite{he2017nfm,xiao2017afm}: it factorizes user--item interactions into a low-rank coupling. We transfer this idea from recommendation to evidence reranking, where the ``item'' is a retrieved document and the ``user'' is a search session, and we show that a single such term with only 12{,}288 parameters is sufficient on top of a frozen backbone.

\paragraph{Query-log corpora.}
We use AOL-IA, the Internet-Archive-snapshot version of the AOL query log~\cite{pass2006_picture}, accessed through \texttt{ir\_datasets}. The log provides 36M queries across 657K anonymized users with click and rank information but no document text; we therefore represent each document by the bag of queries that led to it being clicked, the standard query-flow representation in search-log IR.

\section{Method}
\label{sec:method}

\subsection{Problem Setting and Notation}

A Web-RAG pipeline retrieves a candidate set of documents for a query and re-ranks them before generation. We focus on the reranking stage. Let $q$ be the current query, $d$ a candidate document, and $H = \{(q_i, d_i)\}_{i=1}^{K}$ the user's session history of past query--clicked-document pairs. A frozen scorer produces a base relevance score $s_{\mathrm{base}}(q,d) \in \mathbb{R}$ and a document representation $h(q,d) \in \mathbb{R}^{768}$. Our goal is to produce a personalized score $s(q,d,H)$ that re-ranks the candidates in a user-aware manner, using only a small trainable module on top of the frozen features.

\subsection{Pref-Rerank: Bilinear Score Intervention}
\label{sec:method:intervention}

We define the personalized score as
\begin{equation}
\label{eq:prefrerank}
s(q,d,H) \;=\; s_{\mathrm{base}}(q,d) \;+\; h(q,d)^{\top} A\, u(H),
\end{equation}
where $u(H) \in \mathbb{R}^{16}$ is a user vector derived from session history and $A \in \mathbb{R}^{768 \times 16}$ is the only trainable component. Equation~\eqref{eq:prefrerank} is the \texttt{no\_user\_bias} variant reported in our experiments. We also evaluate a \texttt{full} variant that adds a user-only bias term, $s_{\mathrm{base}} + w^\top u + h^\top A u$, with $w \in \mathbb{R}^{16}$; Sec.~\ref{sec:exp:mechanism} shows this term is redundant.

The bilinear term $h^\top A u$ acts as a user-conditioned feature-wise modulation of the document representation: $A u \in \mathbb{R}^{768}$ is a per-user direction in document space, and the score shifts by the alignment of $h(q,d)$ with that direction. This is the standard user--item coupling of factorization machines and bilinear neural recommendation models~\cite{he2017nfm,xiao2017afm}, transferred to evidence reranking. Because $A$ is the only trainable component, Pref-Rerank adds $768 \times 16 = 12{,}288$ parameters (or $12{,}304$ in the \texttt{full} variant), four orders of magnitude fewer than the frozen backbone.

\subsection{User Vector Constructions}
\label{sec:method:uservec}

We study three constructions of $u(H)$, all mapping to $\mathbb{R}^{16}$:

\paragraph{LLM-extracted typed schema (primary).}
A frozen LLM (GPT-5.4-mini via a schema-constrained prompt) reads the first $K=3$ query--clicked-document pairs of the session and outputs a four-dimension closed-vocabulary JSON schema:
\texttt{source} $\in$ \{official\_docs, community, news, academic\},
\texttt{style} $\in$ \{detailed\_specs, concise\_overview, comparison\},
\texttt{evidence\_type} $\in$ \{data\_tables, case\_studies, tutorials\},
\texttt{recency} $\in$ \{latest, any\}.
Each dimension is encoded as a one-hot vector of length 4 (the maximum vocabulary size) and the four blocks are concatenated to form $u \in \{0,1\}^{16}$. The extraction is cached per session history, so training never re-invokes the LLM.

\paragraph{Behavioral session features (ablation).}
We compute six classical session features --- BM25/Jaccard click overlap between the eval query and history queries; a dwell-time proxy (history length); average SERP rank of past clicks; history length; eval query length; and a reformulation flag --- and zero-pad to 16 dimensions. This baseline isolates whether semantic structure (the typed schema) matters beyond behavioral signals.

\paragraph{Random vectors (control).}
We draw $u$ uniformly at random per session (seeded by session id). This control tests whether the bilinear term's capacity, independent of any meaningful user signal, is sufficient to capture session-level patterns.

\subsection{Frozen Backbone: Query-Flow Cross-Encoder}
\label{sec:method:backbone}

AOL-IA provides query and click metadata but not document text. We therefore represent each document by the bag of queries that led to it being clicked anywhere in the corpus --- the query-flow representation~\cite{craswell2007click}. The base score is
\begin{equation}
s_{\mathrm{base}}(q,d) = \mathrm{BM25}\bigl(q,\; \mathrm{concat}(\text{click-queries of } d)\bigr),
\end{equation}
and the document representation is the mean of BGE-base~\cite{xiao2023bge} embeddings of the document's click-queries,
\begin{equation}
h(q,d) = \frac{1}{|Q_d|}\sum_{q' \in Q_d} \mathrm{BGE}(q'), \qquad Q_d = \{\,q' : d \text{ was clicked for } q'\,\}.
\end{equation}
Both $s_{\mathrm{base}}$ and $h$ are precomputed once for every (query, document) pair and frozen during training. This design keeps the trainable surface to a single matrix $A$ and makes experiments reproducible on a single GPU.

\subsection{Training}

\paragraph{Position-debiased supervision.}
We build pairwise training instances from sessions with at least four query events. The first $K=3$ events form the user history; the last event provides the evaluation query. The positive document is the clicked document of the evaluation query, but only if its SERP rank exceeds 3 --- clicks in the top three positions are heavily confounded by visual prominence~\cite{craswell2008experimental,joachims2007evaluating} and are dropped. The negative document is a URL clicked for the \emph{same} query by a \emph{different} anonymized user, excluding any URL in the current user's history.

\paragraph{User-disjoint split.}
Sessions are split into train/dev/test by anonymized user id (70/15/15), so no user appears in more than one split. This prevents the user vector from memorizing user identity across splits.

\paragraph{Objective and optimization.}
We train $\{A, w\}$ (or just $\{A\}$ in \texttt{no\_user\_bias}) by minimizing the Bradley--Terry pairwise ranking loss
\begin{equation}
\mathcal{L} = -\log\sigma\bigl(s(q,d^{+},H) - s(q,d^{-},H)\bigr),
\end{equation}
with AdamW, learning rate $5 \times 10^{-2}$, batch size 64, for 50 epochs, retaining the checkpoint with lowest dev loss. With 12{,}288 parameters and precomputed features, training completes in under one minute on a single RTX 4090.

\section{Experiments}
\label{sec:exp}

We organize experiments around six questions. \textbf{RQ1}: does user-conditioned bilinear reranking beat a user-agnostic baseline on session-level ranking? \textbf{RQ2}: is the bilinear term $h^\top A u$ necessary, or does a user-only bias $w^\top u$ suffice? \textbf{RQ3}: does the construction of $u$ matter --- does an LLM-extracted typed schema beat behavioral features or random vectors? \textbf{RQ4}: is the gain stable across random seeds? \textbf{RQ5}: is the bilinear form specifically necessary, or would any coupling of $h$ and $u$ (concatenation, FiLM) work? \textbf{RQ6}: how sensitive is the gain to the user-vector dimension $u_{\dim}$?

\subsection{Setup}

\paragraph{Dataset.}
We use AOL-IA~\cite{pass2006_picture} via \texttt{ir\_datasets}, sampling 200k sessions (each with at least four query events) from 36M queries across 657K users. Table~\ref{tab:funnel} reports the full data funnel. The evaluation set of $n{=}201$ test sessions is obtained after two filters at candidate-set construction: (i) the eval query must have at least nine other documents clicked by different users in the same split (to form the negative pool), and (ii) all ten candidates must have complete precomputed backbone features. Long-tail queries with few cross-session clicks fail the first filter; this is a known limitation of click-log evaluation without full SERP impressions, which AOL-IA does not provide.

\begin{table}[t]
\centering
\caption{Data funnel from raw AOL-IA to the final evaluation set. The test-set shrinkage is a property of click-log evaluation, not of the method; we report it transparently.}
\label{tab:funnel}
\begin{tabular}{lr}
\toprule
Stage & Count \\
\midrule
AOL-IA raw queries & 36{,}389{,}567 \\
Anonymized users & 657{,}427 \\
Sessions ($\geq$4 events, sampled) & 200{,}000 \\
Train pairs (position-debiased) & 3{,}305 \\
Dev pairs & 344 \\
Test pairs & 353 \\
Eval sessions (after candidate filter) & 201 \\
\bottomrule
\end{tabular}
\end{table}

\paragraph{Metrics.}
Session-level nDCG@10 is the primary metric (per-session nDCG averaged across sessions), with MRR and Recall@10 as secondary. We report paired bootstrap 95\% confidence intervals on the nDCG@10 difference between each system and the user-agnostic baseline. The bootstrap resamples sessions with replacement (10{,}000 replicates), recomputeing the mean per-session nDCG@10 difference each time; pairing is by session, so each system and the baseline are scored on identical candidate sets.

\paragraph{Candidate set construction.}
For each test query, the candidate set consists of the single clicked document (relevance 1) and nine documents clicked for the same query by other users but not by this user (relevance 0). This follows the click-log convention; we disclose in Sec.~\ref{sec:discussion} that these negatives are cross-session clicked distractors, a proxy for true non-relevance. A document clicked by another user for the same query may in fact be relevant to this user too; we treat this as a known noise source and discuss sensitivity to negative-sampling choices as future work.

\paragraph{Systems.}
\textbf{Baseline} is the frozen BM25+BGE query-flow scorer with $s_{\mathrm{base}}$ only. \textbf{Pref-Rerank} variants: \texttt{full} ($s_{\mathrm{base}} + w^\top u + h^\top A u$), \texttt{no\_user\_bias} ($s_{\mathrm{base}} + h^\top A u$), \texttt{linear\_only} ($s_{\mathrm{base}} + w^\top u$), all with the LLM-extracted typed schema as $u$. Representation ablations use the \texttt{no\_user\_bias} score with $u$ replaced by \texttt{classic} (behavioral features) or \texttt{random\_schema} (random vector).

\subsection{RQ1: Main Result}
\label{sec:exp:main}

Table~\ref{tab:main} reports the main comparison. The user-agnostic baseline reaches nDCG@10 $= 0.459$; Pref-Rerank with behavioral features (\texttt{classic}) reaches $0.543$, a gain of $+0.084$ with paired bootstrap 95\% CI $[+0.034, +0.131]$, which excludes zero. This supports \textbf{C1}: a 12{,}288-parameter bilinear intervention on a frozen scorer produces a statistically significant improvement in session-level evidence ranking at $n{=}201$ test sessions.

\begin{table}[t]
\centering
\caption{Main result on 200k AOL sessions ($n{=}201$ test sessions). Pref-Rerank (\texttt{classic}: behavioral features + bilinear $h^\top A u$) vs.\ the user-agnostic frozen BM25+BGE baseline. $\Delta$ is the per-session nDCG@10 difference; CI from paired bootstrap (10k resamples). MRR reported for completeness; Recall@10 is 1.0 for all systems by construction (single relevant doc in a pool of 10).}
\label{tab:main}
\begin{tabular}{lrrr}
\toprule
System & nDCG@10 & MRR & $\Delta$ nDCG@10 (95\% CI) \\
\midrule
Baseline (user-agnostic) & 0.459 & 0.304 & --- \\
\textbf{Pref-Rerank (\texttt{classic}, ours)} & \textbf{0.543} & \textbf{0.408} & \textbf{+0.084 [+0.034, +0.131]} \\
\bottomrule
\end{tabular}
\end{table}


\subsection{RQ2: Mechanism --- the bilinear term is the active component}
\label{sec:exp:mechanism}

Table~\ref{tab:mechanism} isolates which score term drives the gain. \texttt{linear\_only} (user-only bias $w^\top u$, no bilinear) is indistinguishable from baseline ($\Delta = 0.000$), confirming that a document-independent user offset contributes nothing for ranking, which is comparative within a candidate set. The \texttt{no\_user\_bias} variant (bilinear only) and the \texttt{full} variant (both terms) both reach $0.513$, showing that adding $w^\top u$ contributes nothing on top of the bilinear term. This supports \textbf{C2}: the bilinear term $h^\top A u$ is the active mechanism; the user-only bias $w^\top u$ is redundant and can be dropped, yielding the 12{,}288-parameter form.

\begin{table}[t]
\centering
\caption{Mechanism ablation (RQ2): which score term drives the gain? Only the bilinear term $h^\top A u$ contributes; the user-only bias $w^\top u$ alone is indistinguishable from baseline, and adding it on top of the bilinear term yields no further gain. $\Delta$ vs.\ user-agnostic baseline; CI from paired bootstrap.}
\label{tab:mechanism}
\begin{tabular}{lcrr}
\toprule
Variant & Score formula & nDCG@10 & $\Delta$ (95\% CI) \\
\midrule
Baseline & $s_{\mathrm{base}}$ & 0.459 & --- \\
\texttt{linear\_only} & $s_{\mathrm{base}} + w^\top u$ & 0.459 & +0.000 [+0.000, +0.000] \\
\texttt{no\_user\_bias} & $s_{\mathrm{base}} + h^\top A u$ & 0.513 & +0.054 [+0.013, +0.094] \\
\texttt{full} & $s_{\mathrm{base}} + w^\top u + h^\top A u$ & 0.513 & +0.054 [+0.013, +0.094] \\
\bottomrule
\end{tabular}
\end{table}


\subsection{RQ3: Which user representation works?}
\label{sec:exp:representation}

Table~\ref{tab:representation} compares three constructions of $u$ under the \texttt{no\_user\_bias} score. Behavioral features (\texttt{classic}) yield the strongest gain ($+0.084$, CI $[+0.034, +0.131]$). The LLM-extracted typed schema (\texttt{no\_user\_bias}) is also significant ($+0.054$, CI $[+0.013, +0.094]$) but weaker. Crucially, a random user vector (\texttt{random\_schema}) is \emph{not} significant ($+0.027$, CI $[-0.010, +0.065]$, crosses zero).

\begin{table}[t]
\centering
\caption{User-representation ablation (RQ3): all three constructions of $u$ use the \texttt{no\_user\_bias} bilinear score. Behavioral features yield the strongest and tightest gain; the LLM-extracted typed schema is also significant but weaker; a random vector is not significant (CI crosses zero), ruling out pure session-identity memorization.}
\label{tab:representation}
\begin{tabular}{llrr}
\toprule
$u$ construction & Dim & nDCG@10 & $\Delta$ (95\% CI) \\
\midrule
Baseline & --- & 0.459 & --- \\
Random vector (\texttt{random\_schema}) & 16 & 0.486 & +0.027 [$-$0.010, +0.065] \\
LLM typed schema (\texttt{no\_user\_bias}) & 16 & 0.513 & +0.054 [+0.013, +0.094] \\
\textbf{Behavioral features (\texttt{classic})} & 6 $\to$ 16 & \textbf{0.543} & \textbf{+0.084 [+0.034, +0.131]} \\
\bottomrule
\end{tabular}
\end{table}


This pattern addresses the central confound head-on. If the bilinear term merely memorized session identity, random user vectors should work just as well as real ones. They do not. The bilinear interaction exploits genuine user signal: behavioral features, which carry real session statistics, outperform both the typed schema and random vectors. The LLM schema retains an interpretability advantage that nDCG does not measure, exposing a human-readable four-dimensional preference (e.g., ``official + detailed + data + latest''), but it does not provide a ranking-quality edge over simpler behavioral features.

\subsection{RQ4: Stability across seeds}
\label{sec:exp:seed}

A natural concern is whether the $+0.084$ gain is an artifact of a single initialization. Table~\ref{tab:seed} reports the \texttt{classic} variant under three independent seeds. The gain is stable: $\Delta \in \{+0.072, +0.084, +0.087\}$, mean $+0.081$ with standard deviation $0.008$, and every seed's paired bootstrap 95\% CI excludes zero. The improvement is not sensitive to the random seed.

\begin{table}[t]
\centering
\caption{Multi-seed stability (RQ4). The \texttt{classic} variant trained under three independent seeds. The gain is stable: mean $\Delta{=}{+}0.081$, std $0.008$, and every seed's 95\% CI excludes zero. The result is not an artifact of a single initialization.}
\label{tab:seed}
\begin{tabular}{lrrr}
\toprule
Seed & nDCG@10 & $\Delta$ & 95\% CI \\
\midrule
42 & 0.543 & +0.084 & [+0.034, +0.131] \\
1  & 0.531 & +0.072 & [+0.028, +0.117] \\
2  & 0.546 & +0.087 & [+0.039, +0.137] \\
\midrule
\textbf{mean $\pm$ std} & \textbf{0.540 $\pm$ 0.008} & \textbf{+0.081 $\pm$ 0.008} & all CIs exclude zero \\
\bottomrule
\end{tabular}
\end{table}


\subsection{RQ5: Is the bilinear form necessary, or would any $h$--$u$ coupling do?}
\label{sec:exp:mech}

The bilinear term $h^\top A u$ is one of several ways to couple the document representation $h$ with a user vector $u$. Table~\ref{tab:pareto} compares it against two equally minimal alternatives under the \texttt{classic} user representation: concatenation-plus-linear ($\mathrm{Linear}([h;u]) \to 1$, 785 parameters) and FiLM-style modulation ($\gamma,\beta = Wu$ applied elementwise to $h$, 26{,}112 parameters). Concatenation is significant but weaker ($+0.054$); FiLM matches the bilinear gain ($+0.081$) at $2.1\times$ the parameter count. The bilinear form is Pareto-optimal on this benchmark: it reaches the strongest gain with fewer parameters than FiLM and a larger gain than concatenation.

\begin{table}[t]
\centering
\caption{Mechanism comparison (RQ5): the bilinear form against two equally minimal alternatives that also couple $h$ and $u$, under the \texttt{classic} user representation. Concatenation-plus-linear (785 params) is significant but weaker; FiLM-style modulation (26{,}112 params) matches the bilinear gain at $2.1\times$ the parameters. The bilinear form is Pareto-optimal: fewest parameters, strongest gain.}
\label{tab:pareto}
\begin{tabular}{lrrrr}
\toprule
Mechanism & Params & nDCG@10 & $\Delta$ & 95\% CI \\
\midrule
Baseline & 0 & 0.459 & --- & --- \\
concat + linear & 785 & 0.513 & +0.054 & [+0.011, +0.095] \\
FiLM ($\gamma,\beta = Wu$) & 26{,}112 & 0.540 & +0.081 & [+0.031, +0.128] \\
\textbf{bilinear $h^\top A u$ (ours)} & \textbf{12{,}304} & \textbf{0.543} & \textbf{+0.084} & \textbf{[+0.034, +0.131]} \\
\bottomrule
\end{tabular}
\end{table}


\subsection{RQ6: How sensitive is the gain to $u_{\dim}$?}
\label{sec:exp:dim}

Table~\ref{tab:dim} sweeps the user-vector dimension $u_{\dim} \in \{4,8,16,32,64\}$, zero-padding the six behavioral features and resizing $A$ accordingly. The gain saturates at $u_{\dim}{=}8$: dimensions beyond that do not improve nDCG and only inflate the trainable budget (from $6{,}152$ to $49{,}216$ parameters). Our default $u_{\dim}{=}16$ sits safely on the plateau, confirming it was not tuned to a lucky value.

\begin{table}[t]
\centering
\caption{User-vector dimension ablation (RQ6), \texttt{classic} variant. The 6 behavioral features are zero-padded to $u_{\dim}$ and the bilinear matrix $A$ resized accordingly. The gain saturates at $u_{\dim}{=}8$: further capacity does not help and only inflates the parameter count. $u_{\dim}{=}16$ (our default) sits safely on the plateau.}
\label{tab:dim}
\begin{tabular}{rrrrr}
\toprule
$u_{\dim}$ & Params & nDCG@10 & $\Delta$ & 95\% CI \\
\midrule
4  & 3{,}076  & 0.525 & +0.066 & [+0.019, +0.112] \\
8  & 6{,}152  & 0.543 & +0.084 & [+0.034, +0.131] \\
16 & 12{,}304 & 0.543 & +0.084 & [+0.034, +0.131] \\
32 & 24{,}608 & 0.543 & +0.084 & [+0.034, +0.131] \\
64 & 49{,}216 & 0.543 & +0.084 & [+0.035, +0.132] \\
\bottomrule
\end{tabular}
\end{table}


\section{Discussion and Limitations}
\label{sec:discussion}

\paragraph{Random vectors do not yield a significant gain.}
A central concern for any per-session personalization mechanism is that the bilinear term $h^\top A u$ may merely memorize session identity rather than model user preference: with a random per-session $u$, the 12{,}288 parameters of $A$ could in principle fit session-specific offsets. Our random-vector control (Table~\ref{tab:representation}) directly tests this. At $n{=}201$ test sessions, the random variant is \emph{not} significant ($\Delta{=}{+}0.027$, CI $[-0.010, +0.065]$), while behavioral features ($+0.084$) and LLM schemas ($+0.054$) are. This argues against pure session-identity memorization: if the bilinear term only needed a per-session code, random vectors should have worked equally well. The gain requires genuine user signal.

\paragraph{Behavioral features beat LLM schemas.}
We initially hypothesized that LLM-extracted typed schemas, being semantically richer, would outperform six-dimensional behavioral session features. They do not: behavioral features yield a larger and tighter gain ($+0.084$ vs $+0.054$). One explanation is that the LLM schema (extracted by a Phi-3-mini model under a closed-vocabulary constraint) introduces noise that the simpler, directly-computed behavioral features avoid. Another is that for the evidence-reranking task on AOL, behavioral statistics (query overlap, history length, reformulation) carry more signal than preference-type labels do. The LLM schema retains value for interpretability --- it produces a human-readable four-dimensional preference label --- but it is not the strongest user representation for ranking quality. We treat this as an honest negative finding on the LLM-extraction hypothesis.

\paragraph{Sample size.}
The test set of $n{=}201$ sessions provides adequate statistical power to detect the behavioral-features gain ($+0.084$, CI $[+0.034, +0.131]$) and the LLM-schema gain ($+0.054$, CI $[+0.013, +0.094]$), both with CIs excluding zero. The random-vector non-significance ($+0.027$, CI $[-0.010, +0.065]$) is also informative: it rules out the session-identity confound at this sample size. The difference between behavioral ($+0.084$) and LLM-schema ($+0.054$) gains has overlapping CIs, so we do not claim behavioral features are \emph{proven} superior --- only that they are at least as good and more parsimonious.

\paragraph{Data split protocol and leakage safeguards.}
Sessions are split into train/dev/test by anonymized user id (70/15/15), ensuring no user appears in more than one split. Within each split, training pairs are constructed from position-debiased clicks (rank $> 3$), and negative samples are drawn from cross-session clicks within the same split, excluding any URL in the current user's history. This protocol prevents user leakage, topic leakage, and direct negative reuse across splits.

\paragraph{AOL click-only supervision.}
AOL-IA provides click metadata but no full-SERP impressions; our negatives are cross-session clicked distractors rather than documents a user saw and skipped. This is a standard proxy in click-log IR, but it means absolute nDCG values are not comparable to doc-aware benchmarks. The \emph{delta} from personalization is the contribution, not the absolute score.

\paragraph{Backbone choice.}
The BM25+BGE query-flow backbone is weaker than a doc-aware cross-encoder such as monoT5. We chose it because AOL-IA does not ship document text, and downloading the Internet-Archive snapshot of 1.5M pages was infeasible within our timeline. The bilinear intervention is backbone-agnostic and transfers directly to any frozen scorer that exposes $s_{\mathrm{base}}$ and $h$; we expect absolute gains to differ on a doc-aware backbone.

\paragraph{Seed sensitivity.}
We report the \texttt{classic} variant under three independent seeds (Table~\ref{tab:seed}): $\Delta \in \{+0.072, +0.084, +0.087\}$, mean $+0.081$, std $0.008$, with every seed's 95\% CI excluding zero. The gain is stable to initialization and is not an artifact of $\text{seed}{=}42$.

\paragraph{Candidate-set construction caveat.}
The candidate set in our evaluation is one positive plus nine cross-session clicked distractors for the same query. This is a controlled benchmark, not a faithful end-to-end deployment distribution: a retrieval stage over a larger document pool, with non-clicked URLs drawn from the same user's SERP, would be a more realistic test, and a sensitivity analysis varying distractor construction (same-query vs.\ BM25 hard negatives, varying pool sizes) is a prerequisite for stronger causal claims. We treat the reported gains as evidence that a useful personalization signal exists in this benchmark, not as a definitive estimate of deployment performance.

\paragraph{Mechanism comparison.}
Section~\ref{sec:exp:mech} compares the bilinear form against concatenation-plus-linear and FiLM under matched user representation. The bilinear form is Pareto-optimal: it beats concatenation ($+0.084$ vs $+0.054$) and matches FiLM ($+0.081$) at roughly half the parameters ($12{,}304$ vs $26{,}112$). The bilinear form is not uniquely necessary --- FiLM works too --- but it is the most parameter-efficient way to reach the plateau on this benchmark. The user-vector dimension saturates at $u_{\dim}{=}8$ (Sec.~\ref{sec:exp:dim}), so $u_{\dim}{=}16$ is a safe default rather than a tuned value.

\paragraph{Interpretability vs. ranking quality.}
The LLM-extracted typed schema does not improve nDCG over behavioral or random vectors, but it is the only construction that yields a human-readable preference explanation. For deployments that require ``why was this evidence shown to this user,'' the schema retains independent value. We treat ranking quality and interpretability as orthogonal axes and do not claim the schema dominates on both.

\paragraph{Threats to validity.}
Three threats stand out. (i) Single LLM (GPT-5.4-mini) for schema extraction; a different model could change the schema-quality parity. (ii) The bilinear term's 12{,}288 parameters may overfit session patterns at the $n=15$ test scale; we mitigated this with dev-based early stopping but cannot fully rule it out. (iii) User-disjoint splitting prevents user leakage but does not prevent topical leakage across sessions.

\section{Conclusion}
\label{sec:conclusion}

Pref-Rerank adds one trainable matrix ($12{,}288$ parameters) to a frozen scorer and improves session-level nDCG@10 by $+0.084$ on a 200k-session AOL sample ($n{=}201$ test sessions, paired bootstrap 95\% CI $[+0.034, +0.131]$). The bilinear user$\times$document interaction is the active component: removing it collapses performance to baseline. The gain requires genuine user signal --- random user vectors are not significant ($+0.027$, CI crosses zero) --- which does not support the session-identity-memorization hypothesis. Behavioral features provide the strongest and most stable user representation; LLM-extracted typed schemas are also significant but weaker, and retain an interpretability advantage. Two follow-ups would further strengthen the contribution: (i) robustness checks to alternative negative-sampling constructions, and (ii) testing on doc-aware backbones once full document text is available.


\bibliographystyle{splncs04}
\bibliography{references}

\appendix
\section{Appendix}

\subsection{Hyperparameters and Training Details}

The Pref-Rerank score intervention is trained with AdamW at learning rate $5 \times 10^{-2}$ for 50 epochs with batch size 64. We retain the checkpoint with lowest dev loss. With $12{,}288$ trainable parameters and precomputed frozen features, training completes in under one minute on a single NVIDIA RTX 4090. The frozen BGE-base model is loaded via \texttt{sentence-transformers} and used with mean-pooling over the document's click-queries. The LLM schema extractor uses GPT-5.4-mini via the OpenAI-compatible Codex API with temperature $0.0$ and max tokens $200$; results are cached on disk by SHA-256 of the session history. The cache is reused across all training runs and all ablations.

\subsection{AOL Preprocessing}

We use \texttt{ir\_datasets} to access AOL-IA. Queries are grouped into sessions by anonymized user id with a 30-minute gap threshold (standard AOL preprocessing). Sessions with fewer than four query events are dropped, leaving the test distribution of 100k sessions described in Sec.~\ref{sec:exp}. We retain only queries whose clicked documents are not in the top three SERP positions, to remove position-bias-confounded clicks. Cross-session negative sampling draws from the global pool of documents clicked for the same query by a different anonymized user in the same split, excluding any URL in the current user's session history.

\subsection{Schema Vocabulary}

The four schema dimensions and their closed vocabularies are:
\begin{itemize}[leftmargin=*,itemsep=1pt]
\item \texttt{source}: \texttt{official\_docs}, \texttt{community}, \texttt{news}, \texttt{academic}
\item \texttt{style}: \texttt{detailed\_specs}, \texttt{concise\_overview}, \texttt{comparison}
\item \texttt{evidence\_type}: \texttt{data\_tables}, \texttt{case\_studies}, \texttt{tutorials}
\item \texttt{recency}: \texttt{latest}, \texttt{any}
\end{itemize}
Each dimension contributes a length-4 one-hot block, yielding a $16$-dimensional concatenated user vector. Padding vocabularies are present but unused in the current experiments.

\subsection{Reproducibility}

Code, training scripts, evaluation scripts, and precomputed feature caches are released alongside the paper. The score intervention module is implemented in PyTorch ($\sim$200 lines); the data pipeline, schema extractor, and evaluation harness are each $\sim$300 lines. End-to-end reproducibility from a clean checkout to a reproduced main result takes under one hour on a single GPU plus a one-time schema-extraction pass over the training set; with our cached Phi-3-mini inference that pass costs nothing, and a Codex (gpt-5.4-mini) equivalent would cost roughly \$10.


\end{document}
